% Certonomous CFD Laboratory --- Closure Challenge technical report
% Source of record: demo-output/website/latex/closure_challenge_report.tex
% Every number in this document traces to a repository artifact named in Appendix A
% or to an external publication cited by DOI / arXiv identifier.
\documentclass[11pt,a4paper]{article}

\usepackage[T1]{fontenc}
\usepackage[utf8]{inputenc}
\usepackage[a4paper,margin=25mm]{geometry}
\usepackage{amsmath}
\usepackage{amssymb}
\usepackage{booktabs}
\usepackage{array}
\usepackage{longtable}
\usepackage{graphicx}
\usepackage[protrusion=true,expansion=false]{microtype}
\usepackage{fancyhdr}
\usepackage{xcolor}
\usepackage[hidelinks]{hyperref}

\newcolumntype{L}[1]{>{\raggedright\arraybackslash}p{#1}}
\newcommand{\breakus}{\textunderscore\allowbreak}
\newcommand{\code}[1]{{\small\ttfamily\let\_\breakus #1}}
\newcommand{\artifact}[1]{{\footnotesize\ttfamily\let\_\breakus #1}}

\pagestyle{fancy}
\fancyhf{}
\fancyhead[L]{\footnotesize Certonomous --- Closure Challenge technical report}
\fancyhead[R]{\footnotesize \thepage}
\renewcommand{\headrulewidth}{0.4pt}

\sloppy
\emergencystretch=3em
\hbadness=10000

\setlength{\emergencystretch}{3em}
\tolerance=1200
\hbadness=4000
\sloppy


\title{\textbf{The Closure Challenge}\\[4pt]
\large A post-hoc RANS correction entry, its verification record,\\
and the field-inversion line behind it}
\author{Certonomous CFD Laboratory}
\date{5 August 2026}

\begin{document}
\maketitle
\thispagestyle{fancy}

\begin{center}
\fbox{\begin{minipage}{0.92\textwidth}
\small
\textbf{Status.} Nothing has been submitted to the Closure Challenge. No account has been
created and the benchmark steward has not been contacted. The entry described here
(overall \textbf{0.0654}, round 4) is an internally measured position against the public
leaderboard, not an official ranking. The submission package exists, is audited, and is
one email away; the decision to send it rests with Katie Kim.
\end{minipage}}
\end{center}

\begin{abstract}
\noindent
This report records what the Certonomous CFD laboratory has done in the Closure Challenge,
a public benchmark for machine learning in turbulence modelling stewarded by Ryley McConkey
(MIT). Our entry of record scores \textbf{0.0654} on the benchmark's own unmodified scorer,
which would place it third of five against the current public leaderboard, between
Wu \& Zhang (0.0624) and Liu, Wang, Zhao \& Xiao (0.0737), and 36.9\% below the 0.1036
scored by submitting the supplied RANS fields unchanged. The method is a post-hoc velocity
field correction --- the only entry on the board that does not re-solve the governing
equations with its correction folded in --- combined with a train-only \emph{decline gate}
that withholds the correction on whole cases where it predicts the correction would do
harm, and a pre-registration discipline under which every model choice is frozen in a
committed document before the prediction set it selects is ever scored. We state the costs
of that method class as plainly as its benefits: our submitted fields do not satisfy
continuity, and as of yesterday the departure is measured rather than unmentioned --- about
10\% of the local velocity-gradient scale on the two corrected hills and 2.3--3.4\% on the
ducts, against baselines that are divergence-free to the operator floor; two of the five
leaderboard-best rows we hold are the organisers' own baseline rather than our model; and
one of the remaining three is a lead of 0.00003 that must not be described as a comfortable
win. The report also documents the parallel
field-inversion line, which in the last week moved from a hard numerical blocker
(\code{PetscConvergedReason: -9}, reproduced on two cases and five matrix orderings) to
the laboratory's first converged, finite-difference-verified closure-relevant adjoint
gradient: 21{,}000 design variables on the curved backward-facing step, agreeing with
central differences to 0.021--0.199\% on four independently chosen cells.
\end{abstract}

\vspace{4pt}
\noindent\textbf{Provenance convention.} Every quantitative claim below is followed by the
repository artifact it comes from. Repository paths are relative to the \code{Certonomous}
repository root and are cited with the commit at which the cited text stands; external work
is cited by DOI or arXiv identifier and was read from a copy held under
\artifact{docs/papers/}. Appendix~A is the full provenance index. Where a number is a bound
rather than a measurement, or an inference of ours rather than a claim of the source's, this
report says so at the point of use.

\tableofcontents
\newpage

%=====================================================================
\section{The challenge, and the metric it is scored on}
%=====================================================================

\subsection{What the benchmark is}

The Closure Challenge is an open, ongoing benchmark task for machine learning in turbulence
modelling. Its authoritative surface is the repository
\url{https://github.com/rmcconke/closure-challenge-benchmark}, held locally at commit
\code{deb91557}; the evaluation package \code{closure-challenge} is pinned at commit
\code{1c4e22c8}; the accompanying preprint is arXiv:2603.28884 (McConkey, Buchanan, Smidt,
Bodner, Dwight, Cinnella), DOI \code{10.48550/arXiv.2603.28884}. The README states in its
own words that it, not the paper, is the current source of information, and our records
follow that ordering: no rule in this document is sourced to the preprint where the README
states the same thing
(\artifact{CLOSURE\_CHALLENGE\_SUBMISSION\_DRAFT.md} \S1--\S2, commit
\code{15530f97}).

The task is to predict the velocity field on eight held-out test cases given a suggested
training and validation split and a supplied CFD mesh. The eight scored cases are four
parametric periodic hills (\code{alpha\_15\_13929\_4048}, \code{alpha\_15\_13929\_2024},
\code{alpha\_05\_4071\_4048}, \code{alpha\_05\_4071\_2024}), three rectangular ducts
(\code{AR\_1\_Ret\_360}, \code{AR\_3\_Ret\_360}, \code{AR\_14\_Ret\_180}) and the NASA
wall-mounted hump (\code{NASA\_2DWMH}).

Three properties of the rules matter to everything that follows.

\begin{enumerate}
\item \textbf{There is exactly one strict rule}, and it is about data, not method:
  it is \emph{``strictly forbidden to train or validate on any data from the test cases''},
  on pain of automatic withdrawal and a permanent note on the leaderboard. The preprint
  adds the permission the entry leans on, verified verbatim from the preprint's own
  \S2.1 on 2026-08-02: \emph{``You can train on similar flows to the test cases.''}
\item \textbf{Everything else is unrestricted.} Input features, base turbulence model and
  data-assimilation technique are explicitly the submitter's choice; the published
  train/validation split is labelled \emph{``(suggested)''} and only the test column binds.
\item \textbf{There is no scoring-call limit, and there cannot be one.} The evaluation
  package ships the test ground truth locally and the README instructs submitters to
  preview their score; \code{evaluate\_by\_case()} is documented public API. Our
  five-calls-ever ledger (\S\ref{sec:ledger}) is therefore a self-imposed discipline,
  \emph{stricter} than the benchmark requires, and must never be described as compliance
  with a rule that does not exist.
\end{enumerate}

There is no eligibility clause of any kind in the benchmark README, the evaluation package
README or the preprint --- an absence of restriction, reported here as an absence rather
than as a permission. No fee, registration, account or affiliation is required; submission
is an email of a \code{test/} directory to the steward.

\subsection{The metric, read from the evaluation code}

The metric was read directly from
\code{closure\_challenge/eval.py} rather than from our own prose
(\artifact{CLOSURE\_METHODS\_COMPARISON.md} \S1.0, commit
\code{ac2f37ee}). For a case with 1000 fixed evaluation points,
\code{evaluate\_individual\_case()} computes a scaled mean absolute error of the
\emph{vector magnitude} of the velocity error,
%
\begin{equation}
s_{\text{case}} \;=\;
\frac{\dfrac{1}{N}\sum_{i=1}^{N}\bigl\lVert \mathbf{U}^{\text{pred}}_i - \mathbf{U}^{\text{true}}_i \bigr\rVert_2}
     {\dfrac{1}{N}\sum_{i=1}^{N}\bigl\lVert \mathbf{U}^{\text{true}}_i \bigr\rVert_2},
\qquad N = 1000,
\end{equation}
%
and \code{score()} returns the plain, unweighted mean of the eight per-case values. Lower
is better. Two consequences are worth carrying: the error is a vector-magnitude error and
not component-wise, and because the mean is unweighted, an improvement of $\Delta$ on one
case moves the overall by $\Delta/8$ --- so the two \code{alpha\_15} cases, whose baseline
errors of 0.1320 and 0.2049 dominate the floor, carry exactly the same $1/8$ weight as
everything else.

All scoring in this project is done through the benchmark's own unmodified scorer. The
metric is never reimplemented for scoring purposes.

\subsection{The public leaderboard, and where we stand}

Table~\ref{tab:board} gives the public board as published at benchmark commit
\code{deb91557}, with our unsubmitted entry inserted at its measured position. All four
published entrants re-score to their published values exactly on our harness, per-case and
overall, which is the strongest available evidence that the harness is being driven
correctly (\artifact{CLOSURE\_CHALLENGE\_SUBMISSION\_DRAFT.md} \S7.5).

\begin{table}[h]
\centering
\small
\caption{The public leaderboard, with our entry of record at its internally measured
position. We are unsubmitted; this is not an official ranking.}
\label{tab:board}
\begin{tabular}{clc}
\toprule
Rank & Entry & Overall \\
\midrule
1 & Reissmann, Fang \& Sandberg & 0.0595 \\
2 & Wu \& Zhang & 0.0624 \\
--- & \textbf{Certonomous (unsubmitted, round 4, entry of record)} & \textbf{0.0654} \\
--- & Certonomous (unsubmitted, round 3, superseded) & 0.0676 \\
3 & Liu, Wang, Zhao \& Xiao & 0.0737 \\
4 & Montoya, Oulghelou \& Cinnella & 0.0779 \\
\midrule
--- & RANS-identity floor (submit the supplied fields unchanged) & 0.1036 \\
\bottomrule
\end{tabular}
\end{table}

Two cautions attach to that table and are repeated wherever it appears.

\begin{itemize}
\item \textbf{The 0.1036 floor is our own number.} The script that produces it
  (\code{scripts/rans\_identity\_baseline.py} in the benchmark clone) is untracked --- it
  was written by this laboratory, not shipped by the benchmark, and the README quotes no
  RANS-identity baseline anywhere. It reproduces exactly and is honestly derived, but any
  wording implying the floor is the benchmark's own published figure is false
  (\artifact{CLOSURE\_CHALLENGE\_SUBMISSION\_DRAFT.md} \S7.2).
\item \textbf{Position claims go stale.} The challenge is ongoing with no deadline; the
  board has been unchanged since 2026-05-04 but can change without notice. Every position
  claim in this document is dated 5 August 2026.
\end{itemize}

%=====================================================================
\section{Our approach}
%=====================================================================

\subsection{Method class, stated at its real strength}
\label{sec:method}

Our correction is a \textbf{post-hoc velocity-field correction}. The learning target is
$\delta \mathbf{U} = \mathbf{U}_{\text{LES}} - \mathbf{U}_{\text{RANS}}$ per mesh cell, and
the learned map is applied once to an already-converged RANS field as a final
post-processing step. \textbf{Nothing re-enters the governing equations and nothing is
re-solved with the correction folded in.} In the Duraisamy--Iaccarino--Xiao taxonomy ---
\emph{Turbulence Modeling in the Age of Data}, \emph{Annu.\ Rev.\ Fluid Mech.} 51 (2019)
357--377, cited at its honest tier, metadata only: the review has not been read in full by
this laboratory, so ``correct the answer'' is our own paraphrase of the family via
\artifact{docs/research/CLOSURE\_METHODS.md}, not the review's wording --- this is
structurally distinct from FIML, TBNN, SpaRTA and
eigenvalue perturbation, all of which correct something inside the equations before
re-solving (\artifact{CLOSURE\_CHALLENGE\_STATUS.md} \S4, commit \code{b9cb56a2}).

\paragraph{Correction, 2026-08-05 --- the family we are in now carries its citation.} The
passage above previously placed the method in a taxonomy, named four families we are not,
and cited no prior work for the family we are in --- origination by omission, named as such
by the method-priority review (\artifact{CLOSURE\_METHOD\_PRIORITY\_REVIEW.md} \S4.1,
commit \code{d84b649f}). The owed citation is paid here, at read-in-full tier: learned
local-error correction of a cheap CFD solve, applied once without re-solving, is
established --- Hanna, Dinh, Youngblood \& Bolotnov, \emph{Coarse-Grid Computational Fluid
Dynamic (CG-CFD) Error Prediction using Machine Learning}, arXiv:1710.09105 (2017);
journal version \emph{Prog.\ Nucl.\ Energy} 118 (2019) 103140, DOI
10.1016/j.pnucene.2019.103140 (journal record: metadata only). The arXiv full text was
fetched and read on 2026-08-05 and is held at
\artifact{docs/papers/hanna\_dinh\_youngblood\_bolotnov\_1710.09105.pdf}. Their surrogate
predicts the cheap solve's local error from the cheap solve's own local features and adds
it back to the variables of interest, nothing re-solved --- our architecture with the
error source swapped: their cheap solve is a coarse-grid \emph{no-model} Navier--Stokes
solve of a lid-driven cubic cavity and the corrected error is grid-coarsening
(discretisation) error, tested across unseen Reynolds numbers and grid sizes; ours is a
converged $k$--$\omega$ SST solve and the error is turbulence-closure error on separated
and secondary flows. Two deltas from the full read are worth carrying beyond the review's
abstract-tier characterisation: their 37 features are the cell Reynolds number plus scaled
first and second velocity derivatives --- grid-resolution features, not physics
invariants --- and their own open-issues section names the exact cost our G2 audit later
measured on ourselves, that the velocity components ``are corrected separately, without
enforcing conservation and Galilean invariance''. The founding CFD paper of the class
states the class's physicality gap; it does not soften our measured instance of it
(\S\ref{sec:stability}). The statistical ancestor is model-discrepancy calibration ---
Kennedy \& O'Hagan, \emph{J.\ R.\ Statist.\ Soc.\ B} 63(3) (2001) 425--464, cited as
lineage at metadata-only tier: not read in full by this laboratory, and no sentence of it
is quoted. What survives, stated in the same breath: within the data-driven RANS-closure
literature swept by the review (27 recorded searches, 13 counted negatives), every located
correction that produces a velocity field re-enters the equations and re-solves; no other
post-hoc \emph{velocity-field} correction was found, so the application to turbulence
closure is ours even though the class is not.

\paragraph{The published argument for the target choice.} One prior argues \emph{for}
this method class rather than merely tolerating it: Wu, Xiao, Sun \& Wang,
\emph{Reynolds-averaged Navier--Stokes equations with explicit data-driven Reynolds
stress closure can be ill-conditioned}, \emph{J.\ Fluid Mech.} 869 (2019) 553--586,
arXiv:1803.05581 (abstract read 2026-08-05). Propagating a learned Reynolds-stress
correction through the RANS equations is a published conditioning hazard --- per the
arXiv record's own summary, stress errors below 0.5\% can amplify to velocity errors up
to 35\%; those figures are quoted at the record's tier and were not read off the paper's
tables --- and every re-solving entrant on this board must survive it, where correcting
the velocity field directly sidesteps it. This citation defends the choice of target as
principled rather than merely cheap; it says nothing about the continuity cost measured
in \S\ref{sec:stability} and must not be used to soften it.

Two model objects stand behind the six corrected predictions of the entry of record
(\artifact{CLOSURE\_METHODS\_COMPARISON.md} \S1.1):

\begin{itemize}
\item \textbf{The periodic-hill (PH) model} --- three \code{HistGradientBoostingRegressor}s,
  one per velocity component, \code{max\_iter=300}, \code{max\_depth=6},
  \code{learning\_rate=0.05}, \code{l2\_regularization=1.0}, \code{random\_state=0}, on
  seven per-cell features: Pope's five scalar invariants of the normalised strain and
  rotation tensors (\code{I1\_S2}, \code{I2\_W2}, \code{I3\_S3}, \code{I4\_W2S},
  \code{I5\_W2S2}; Pope, \emph{J.\ Fluid Mech.} 72 (1975) 331--340), a wall-distance
  turbulent Reynolds number \code{Re\_y}, and a turbulent-to-mean kinetic energy ratio
  \code{tke\_ratio}. Trained on 21 periodic-hill cases, 327{,}600 cells.
\item \textbf{The duct model, ``Variant D''} --- the same architecture on eight features
  (adding a Reynolds-invariant $d/d_{\max}$) with a velocity-scale-normalised target,
  $(\mathbf{U}_{\text{LES}} - \mathbf{U}_{\text{RANS}})/\overline{|\mathbf{U}_{\text{RANS}}|}$
  per case, trained on the four suggested training ducts, 41{,}971 cells.
\end{itemize}

Round 2 required building a feature pipeline the benchmark does not ship for the duct and
CBFS families: Green--Gauss $\nabla\mathbf{U}$ reconstructed from mesh, RANS $\mathbf{U}$
and boundary conditions, and wall distance from the nearest wall-patch face centre. It was
validated \emph{only} against shipped fields on periodic-hill \emph{training} cases: cell
centres and volumes to $\sim$1e-13 and $\sim$1e-22 absolute, $\nabla\mathbf{U}$ to a pooled
relative Frobenius error of 0.0012--0.0029 with Pearson $r = 0.9997$--$1.0000$ on three
different training cases, and wall distance to $r = 1.0000$ with mean relative error below
0.1\% (\artifact{CLOSURE\_CHALLENGE\_STATUS.md} \S4).

\subsection{The do-no-harm gate}
\label{sec:gate}

The distinguishing component of the entry is a \textbf{case-level decline gate}. It is a
three-feature ridge regression (top-3 screened features \code{p90\_I4\_W2S},
\code{frac\_backflow}, \code{p90\_I3\_S3}; RidgeCV $\alpha = 0.7499$) that predicts, from
RANS-derived quantities alone, the scaled MAE the \emph{uncorrected baseline} will achieve
on a case. If the predicted baseline error falls below a threshold of 0.1263 --- the
leave-one-out median over the training cases --- the correction is withheld and the
unmodified RANS field \emph{is} the submission for that case
(\artifact{closure\_challenge\_C1\_error\_estimator.md}, commit
\code{5719374e}).

Every parameter of the gate --- feature screening, ridge $\alpha$, standardisation, final
coefficients and threshold --- was computed from the 21 periodic-hill training cases and
nothing else, and was validated on the four suggested validation cases (AUC 1.0, 4 of 4
correct) before it was ever pointed at a test case. When it was applied, the refit
reproduced $\alpha = 0.7499$ and threshold $0.1263$ byte-for-byte from train-only data as a
precondition of the run.

Two honest qualifications belong beside it:

\begin{itemize}
\item \textbf{Its statistical base is thin, and this was disclosed before it was
  useful.} AUC 1.0 on $n=4$ with a single positive is approximately a one-in-four
  outcome by luck; the gate has made eight binary decisions in total. Widening the
  validation base to a 25-case leave-one-out is a filed proposal (G5,
  \S\ref{sec:audit}).
\item \textbf{It is not novel, and the nearest prior art is by the challenge's own
  authors.} Two established things, and they are not the same thing --- the list below
  was one undifferentiated group until 2026-08-05, when the method-priority review
  (\artifact{CLOSURE\_METHOD\_PRIORITY\_REVIEW.md} \S4.2, commit \code{d84b649f}) found
  that grouping credited two of the four papers with a mechanism they did not report.
  \textbf{Identifying} where the baseline is unreliable, from RANS-only inputs, is
  established: Ling \&
  Templeton, \emph{Phys.\ Fluids} 27, 085103 (2015) --- whose abstract, read at the OSTI
  record, classifies RANS results point by point as high or low uncertainty and controls
  no correction --- and Wu, Wang, Xiao \& Ling,
  \emph{Flow Turbul.\ Combust.} 99, 25 (2017), DOI 10.1007/s10494-017-9807-0, an
  \emph{a priori} confidence measure. \textbf{Using such a classifier to control where a
  data-driven correction is fitted and applied} is established separately: Steiner,
  Dwight \& Vir\'e, \emph{Flow Turbul.\ Combust.} (2022), DOI 10.1007/s10494-022-00346-6;
  and Buchanan, L\u{a}c\u{a}tu\c{s}, West \& Dwight, \emph{Computers and Fluids} 305,
  106899 (2025), DOI 10.1016/j.compfluid.2025.106899, arXiv:2504.06758 --- whose first
  and last authors are co-authors of the Closure Challenge paper itself. Read in full,
  their RITA classifier selects \emph{regions inside a case} from local ratios of the
  baseline's own $k$-equation terms. What survives of our distinction, and all we claim,
  is: our gate acts on a \emph{whole case}; it predicts the \emph{baseline's error}
  rather than distribution shift; and its ``off'' state means the uncorrected field
  \emph{is} the submission
  (\artifact{CLOSURE\_CHALLENGE\_PRIOR\_ART.md} \S2.4, \S7.4).
\item \textbf{The nearest miss, quoted rather than characterised} \emph{(added
  2026-08-05, method-priority review \S2.2)}. Ji, Luo, Zhou \& Zhao, \emph{Progressive
  Mixture-of-Experts with autoencoder routing for continual RANS turbulence modelling},
  arXiv:2601.09305 (2026), \textbf{read in full} 2026-08-05, computes exactly our
  quantity --- a whole-flow confidence from the baseline field alone, ``the relative
  frequency of recognised points'' --- and then does something else with it: ``if
  $p_\mathcal{K}$ falls below $T_{\text{accept}}$, the flow is deemed unknown, which
  triggers the continual learning process.'' Its low-confidence branch \emph{expands},
  spawning an expert. There is no branch that submits the uncorrected baseline. That is
  the single closest published approach to our gate and it stops one step short --- which
  is the strongest available support for our narrow claim, and it is narrow.
\end{itemize}

That last reading carried a constraint we wrote down before anyone felt the pull of it:
Buchanan et al.\ train on the NASA wall-mounted hump, which is one of our eight scored test
cases. Their published model coefficients are therefore \textbf{off limits} to this entry
--- borrowing them or calibrating anything of ours against that model's hump behaviour
would make our submission indirectly trained on a test case.

\subsection{Pre-registration}

Every model choice that selects between candidate prediction sets is frozen in a committed
document \emph{before} the prediction set it selects is scored. The NASA hump's model
choice was pre-registered before the round-2 scoring call; Variant D was pre-registered on
the benchmark's own suggested duct validation case \code{AR\_7\_Ret\_180} before any test
score for it existed
(\artifact{closure\_challenge\_duct\_reynolds\_transfer.json}, commit
\code{d0992da2}); the round-5 candidate wrote down its accept/reject rule, then
\emph{failed its own pre-declared gate and was not scored} (\S\ref{sec:nogo}).

Pre-registration that only ever ratifies is decoration. The R5 no-go is the evidence that
this one binds.

\paragraph{Whose discipline this is, added 2026-08-05.} Neither half of this practice is
ours, and the method-priority review found this the element our records named the fewest
priors for (\artifact{CLOSURE\_METHOD\_PRIORITY\_REVIEW.md} \S2.3, commit
\code{d84b649f}). Freezing an analysis plan before results is an established practice in
machine learning with its own venue: the \textbf{NeurIPS 2020 Workshop on
Pre-registration in Machine Learning}, PMLR volume 148, held 11 December 2020 (editors
Bertinetto, Henriques, Albanie, Paganini, Varol), with a second edition as PMLR volume
181 (NeurIPS 2021) --- cited at metadata-only tier from the publisher's volume pages; no
paper from either volume has been read here. And the problem our scoring-call ledger
(\S\ref{sec:ledger}) addresses has a canonical paper: Blum \& Hardt, \emph{The Ladder: A
Reliable Leaderboard for Machine Learning Competitions}, arXiv:1502.04585, ICML 2015,
PMLR 37 (abstract read 2026-08-05). Their abstract prices our answer for us: repeated
evaluation against a holdout lets participants overfit it, and existing approaches
``resort to poorly understood heuristics such as limiting the bit precision of answers
and the rate of re-submission.'' \textbf{Our five-call ledger is rate limiting --- one of
those heuristics, self-imposed --- and it should be read as harm reduction rather than a
guarantee.} Blum \& Hardt build a principled alternative and we do not implement it. What
we did not find in that literature is an entrant imposing either discipline on itself,
unasked, on a benchmark that imposes none, and publishing the ledger; that is a
disclosure practice, not a method, and it is claimed as one.

\subsection{The costs of this method class, stated before the results}

Table~\ref{tab:dimensions} compares the five entries dimension by dimension, from each
entrant's own documents (\artifact{CLOSURE\_METHODS\_COMPARISON.md} \S3.1). Read the top
two rows first, because they are the ones that cost us.

\begin{table}[h]
\centering
\footnotesize
\caption{The five entries, dimension by dimension. ``Absent from their documents'' means
absent from the sources read --- two submission documents, one notebook plus info file, and
two full papers --- and is stated as an absence in those sources, not as a claim about any
laboratory's unpublished practice.}
\label{tab:dimensions}
\begin{tabular}{L{2.3cm}L{2.5cm}L{2.1cm}L{2.1cm}L{2.1cm}L{2.1cm}}
\toprule
Dimension & \textbf{Certonomous} & Reissmann & Wu \& Zhang & Liu et al. & Montoya et al. \\
\midrule
Corrected object &
\textbf{velocity field, post-hoc} &
anisotropy + SST coefficients &
$\beta$ on $\omega$-destruction + fixed QCR &
TBNN anisotropy $+$ $k$--$\omega$ coefficients, coupled &
blended tensor-basis + $R_k$ \\
\addlinespace
Re-solved with the correction &
\textbf{no --- unique among the five} &
yes (inferred) & yes & yes & yes \\
\addlinespace
Conservation satisfied by the submitted field &
\textbf{no --- and now measured} (\S\ref{sec:stability}) &
yes, by construction & yes & yes & yes \\
\addlinespace
Trains on a scored case &
no --- enforced by executable asserts &
no & no &
\textbf{yes, per their own paper} (entry-level status unverifiable) &
no \\
\addlinespace
Case-level do-no-harm gate &
\textbf{yes} &
absent from their documents &
region shield $f_d$, hand-designed &
absent from their documents &
baseline-as-expert local fallback \\
\addlinespace
Pre-registration before scoring &
\textbf{yes} & absent & absent & absent & absent \\
\addlinespace
Scoring-call ledger &
\textbf{yes, 5 calls, self-imposed} &
none visible & none visible & none visible & none visible \\
\bottomrule
\end{tabular}
\end{table}

The first three rows are the honest headline of this report. \textbf{Every entrant above us
ships a method whose output satisfies the governing equations by construction; ours does
not.} The submitted corrected field is a cell-wise machine-learning output added to a
solenoidal RANS solution, and it does not satisfy continuity by construction. No
realizability question arises --- we predict velocity, not Reynolds stresses, so that
column is not-applicable rather than passed --- but $\nabla\cdot\mathbf{U}$ is a real
property of a submitted field and until this week nobody had measured it. It has now been
measured against a threshold fixed before the run, and the finding is material and against
us (\S\ref{sec:stability}).

The compensating claim is narrow and we keep it narrow: among the five entries on this
board, ours is the only one that will decline to correct a whole case and submit the
baseline instead, and the only one that freezes its choices in advance of scoring.

%=====================================================================
\section{Results, round by round}
%=====================================================================

\subsection{The arc}

\begin{description}
\item[Round 1 --- 0.0869.] The PH model, trained on 21 periodic-hill cases and validated on
  4, applied to the four periodic-hill test cases only. All four non-PH test cases were
  passed through as unmodified raw RANS, identical to the floor, because the benchmark does
  not ship $\nabla\mathbf{U}$ and wall distance in solved time directories for the duct and
  CBFS families and reconstructing them was out of scope for that round.
\item[Round 2 --- 0.0741 ($\Delta$ $-$0.0128).] The mesh-reconstruction pipeline of
  \S\ref{sec:method} unlocked the remaining four cases. A new duct-specific model was
  trained on \code{AR\_1/3/5/10\_Ret\_180} and validated on the held-out
  \code{AR\_7\_Ret\_180}; \code{NASA\_2DWMH} was corrected with the existing PH model, a
  pre-registered choice made before test scoring. The hump regressed slightly, from 0.0621
  to 0.0632, and was reported rather than reverted.
\item[Round 3 --- 0.0676 ($\Delta$ $-$0.0065).] The C1 decline gate, built and validated in
  a prior session on non-test data only, was evaluated on the four periodic-hill test
  cases' RANS-derived features for the first time. It declined both \code{alpha\_05} cases
  (predicted baseline 0.0420 and 0.0034, both below the 0.1263 threshold) and kept the
  correction on both \code{alpha\_15} cases (0.1503 and 0.1525, both above) --- 4 of 4
  correct on cases it had never been evaluated on, matching its validation performance
  exactly. The movement, $-$0.0065, matched the $-$0.0066 predicted by the round-2 error
  decomposition almost exactly. Compute: 83\,s wall on a two-core cap, peak RSS 356\,MB.
\item[Round 4 --- 0.0654 ($\Delta$ $-$0.0022), the entry of record.] A duct-only change.
  Variant D replaced the round-2/3 duct model on all three ducts. The five non-duct
  predictions are \emph{byte-identical} to round 3 --- copied, not regenerated, and
  SHA-256-verified against the entry manifest --- so the entire delta is attributable to
  the duct family and the gate's behaviour is untouched.
\end{description}

Round 4's diagnosis was pre-registered and its sign pattern was predicted before it was
measured. The published duct diagnosis had been \emph{wrong}: \code{I3\_S3} and
\code{I4\_W2S} are indeed algebraically zero on every duct, but they are equally zero on
\code{AR\_14\_Ret\_180}, the duct we score best on, so they discriminate nothing. What
discriminates is that \code{Re\_y} reaches $1.85\times$ and $2.07\times$ its trained
maximum on the two ducts we were losing, against $0.90\times$ on the one we win --- and a
gradient-boosted tree cannot extrapolate. The pre-registration predicted improvement
concentrated on the two extrapolating cases and little or none elsewhere. Measured:
$-$0.0108, $-$0.0087, and $+$0.0022. The sign pattern matches on all three
(\artifact{CLOSURE\_CHALLENGE\_STATUS.md} \S0e).

\subsection{The per-case ledger}

\begin{table}[h]
\centering
\footnotesize
\caption{Per-case scaled MAE across rounds, against the RANS-identity floor, rank 2, and
the best published score on each case. Lower is better. Round-4 values are the entry of
record.}
\label{tab:percase}
\resizebox{\textwidth}{!}{%
\begin{tabular}{lccccccl}
\toprule
Case & Floor & R2 & R3 & \textbf{R4} & Rank 2 & Best & Attribution \\
\midrule
\code{alpha\_15\_13929\_4048} & 0.1320 & 0.0501 & 0.0501 & \textbf{0.0501} & 0.0813 & 0.0592 & \textbf{our model leads} \\
\code{alpha\_15\_13929\_2024} & 0.2049 & 0.1011 & 0.1011 & \textbf{0.1011} & 0.1195 & 0.1195 & \textbf{our model leads} \\
\code{alpha\_05\_4071\_4048} & 0.0461 & 0.0723 & 0.0461 & \textbf{0.0461} & 0.0569 & 0.0569 & \emph{baseline, declined} \\
\code{alpha\_05\_4071\_2024} & 0.0719 & 0.0974 & 0.0719 & \textbf{0.0719} & 0.0848 & 0.0760 & \emph{baseline, declined} \\
\code{AR\_1\_Ret\_360}       & 0.1288 & 0.0919 & 0.0919 & \textbf{0.0811} & 0.0455 & 0.0387 & behind, 3rd of 5 \\
\code{AR\_3\_Ret\_360}       & 0.1243 & 0.0862 & 0.0862 & \textbf{0.0775} & 0.0399 & 0.0341 & behind, 3rd of 5 \\
\code{AR\_14\_Ret\_180}      & 0.0590 & 0.0303 & 0.0303 & \textbf{0.0325} & 0.0350 & 0.0325 & lead of 0.00003 only \\
\code{NASA\_2DWMH}           & 0.0621 & 0.0632 & 0.0632 & \textbf{0.0632} & 0.0364 & 0.0364 & behind, last on board \\
\midrule
\textbf{Overall} & \textbf{0.1036} & \textbf{0.0741} & \textbf{0.0676} & \textbf{0.0654} & \textbf{0.0624} & --- & rank 3 of 5 \\
\bottomrule
\end{tabular}}
\end{table}

\subsection{How the ``best on board'' count must be reported}
\label{sec:attribution}

Arithmetically we hold the best score on the public board on five of eight cases. That
count was on our public page as a headline and \textbf{we took it down ourselves}, because
it is true arithmetic and misleading attribution
(\artifact{CLOSURE\_CHALLENGE\_SUBMISSION\_DRAFT.md} \S8.2). On
\code{alpha\_05\_4071\_4048} and \code{alpha\_05\_4071\_2024} our submitted field \emph{is}
the organisers' own unmodified RANS solve --- re-derived independently from the supplied
$k$--$\omega$ SST field and mesh and diffed against the shipped CSV to a maximum deviation
of $5.0\times10^{-10}$, which is the CSV's own write precision. Those two values are
properties of a file every entrant is handed for free; anyone submitting it unchanged
scores exactly the same. The honest attribution is Table~\ref{tab:attribution}.

\begin{table}[h]
\centering
\small
\caption{Attribution of the eight test cases, replacing the retired ``5 of 8'' headline.}
\label{tab:attribution}
\begin{tabular}{L{5.0cm}cL{7.2cm}}
\toprule
Attribution & Count & Cases \\
\midrule
Our model, clear lead & \textbf{2} & \code{alpha\_15\_13929\_4048}, \code{alpha\_15\_13929\_2024} \\
Our model, nominal lead only & 1 & \code{AR\_14\_Ret\_180} \\
The baseline led; the gate withheld our model & \textbf{2} & \code{alpha\_05\_4071\_4048}, \code{alpha\_05\_4071\_2024} \\
Behind & 3 & \code{AR\_1\_Ret\_360}, \code{AR\_3\_Ret\_360}, \code{NASA\_2DWMH} \\
\bottomrule
\end{tabular}
\end{table}

What \emph{is} defensible is a better finding than the one it replaces, and it is about the
selection rather than the score: crossing the recorded per-case floor against the published
per-case leaderboard shows that \textbf{on exactly two of the eight test cases the
uncorrected baseline beats all four published entries --- and they are exactly the two
cases the gate declined.} A rule fitted on 21 training cases and validated on 4 non-test
cases, having never seen a test case, identified blind the two flows on which every
published method in this field damages the answer. That is a finding about the benchmark
which the steward may value more than our own number, and the draft cover email says so.

\subsection{The \code{AR\_14\_Ret\_180} regression, deliberately not reverted}

Round 4 improved two ducts and made a third worse: \code{AR\_14\_Ret\_180} went from 0.0303
to 0.0325, $+$0.0022. It was not reverted, and the reason is the same reason the NASA
hump's $+$0.0011 was not reverted in round 2: choosing per case between two models
\emph{after} seeing their per-case test scores is selection on test outcomes, which is
exactly what the benchmark's one strict rule exists to prevent. Variant D was frozen on
validation evidence and applied to all three ducts or none. It was applied to all three.

\medskip
\noindent\fbox{\begin{minipage}{0.95\textwidth}
\small
\textbf{Margin warning, stated so it cannot be misquoted.} \code{AR\_14\_Ret\_180}'s
best-on-board lead over Reissmann collapsed from 0.0022 to \textbf{0.00003} (ours
0.0324698, against the published four-decimal 0.0325). \textbf{That is a nominal lead, not
a meaningful one, and it must not be reported as a comfortable win.} The same case carries
the entry's weakest numerical-provenance column, which makes any unmeasured numerical slop
on it material.
\end{minipage}}

\subsection{Three refusals}

Three separate score improvements were identified, priced, and declined --- each because
taking it would have required choosing what to submit case by case from the answer key.

\begin{enumerate}
\item \textbf{0.0066, refused in round 2.} Switching the correction off on the three cases
  where round 2 was worse than doing nothing would have scored 0.0675. Knowing \emph{which}
  three required reading round-2's per-case test scores. The honest route --- build a
  train-only gate first, validate it on non-test data, then apply it --- recovered 0.0065
  of the same term in round 3 and cost 0.0001. The laboratory's own compliance audit calls
  this ``the most defensible decision in the record''.
\item \textbf{The \code{AR\_14\_Ret\_180} revert, refused in round 4} --- worth 0.0022 on
  one case, 0.000275 on the mean.
\item \textbf{0.00014, refused in 2026-08-01.} On \code{NASA\_2DWMH} our correction is
  worse than not correcting (0.0632 against a floor of 0.0621); submitting the untouched
  solve there is one file copy, breaks no benchmark rule, and would move the overall from
  0.06543 to 0.06529. The only reason to single out that one case is that the benchmark
  told us its score. It is the smallest of the three refusals and it was refused on
  identical grounds (\artifact{CLOSURE\_CHALLENGE\_SUBMISSION\_DRAFT.md} \S8.3).
\end{enumerate}

%=====================================================================
\section{Where the deficit lives, and what separates us from rank 2}
%=====================================================================

\subsection{The decomposition, and the mechanism it found}

The C2 record decomposes the round-2 result case by case
(\artifact{closure\_challenge\_C2\_error\_decomposition.md}, commit
\code{1da013a8}). Its sanity check ran first: recomputing the eight-case means from the
per-case values reproduces the published floor 0.1036 and entry 0.0741 exactly.

Its headline was uncomfortable and is the reason the gate exists: \textbf{on three of eight
cases our correction was worse than doing nothing}, and the two \code{alpha\_05}
regressions ($+$0.0262 and $+$0.0255) were an order of magnitude larger than the already
known NASA hump regression ($+$0.0011), accounting for 98\% of the 0.0528 of summed damage.

The first explanation offered --- an ``alpha-regime failure'' --- was replaced by the same
document. Two facts kill it: the training split already contains five \code{alpha\_05}
cases, with unremarkable per-case training MAE of 0.0589--0.0719, so the model was not
extrapolating into an unseen regime; and \code{NASA\_2DWMH} is not a periodic hill at all
yet sits in the same degraded group. The framing that survives is
\textbf{the correction hurts where RANS was already good}: ordering the five PH-model cases
by baseline quality separates them without overlap, hurting at every floor $\leq 0.0719$
and helping at every floor $\geq 0.1320$, a gap of 0.0601 with nothing in it, Pearson
$r = -0.9418$. \emph{Caveat kept from the source: $n = 5$. A correlation of $-0.94$ on five
points is suggestive, not established.} What earns it precedence over the alpha framing is
not the $r$ value but that it explains all three regressions with one mechanism, including
the one the alpha story cannot touch.

\subsection{Self-correction as a working practice}

C2 is the most-corrected document in the closure line, and this report presents that as a
feature rather than an embarrassment. Table~\ref{tab:corrections} lists its corrections in
place, each dated, each left visible in the record beside the text it corrects rather than
silently rewritten.

\begin{table}[h]
\centering
\footnotesize
\caption{The C2 record's corrections, all made in place and left visible. The docket
describes C2 as ``thrice-corrected''.}
\label{tab:corrections}
\begin{tabular}{L{2.1cm}L{4.6cm}L{7.4cm}}
\toprule
Date & What was corrected & What replaced it \\
\midrule
in document &
The ``alpha-regime failure'' framing &
Baseline quality: the correction hurts where RANS was already good. Retained with its
$n=5$ caveat. \\
\addlinespace
addendum &
\textbf{The floor was the wrong yardstick.} Measuring against ``do nothing'' made the ducts
look like wins ($-$28.6\%, $-$30.7\%) &
Measured against the published board, \code{AR\_1} and \code{AR\_3} together are
\textbf{62.9\%} of the gap to rank 2, and the two \code{alpha\_05} cases are the
\emph{smallest} recoverable terms (18.9\% combined). The earlier recommendation to target
\code{alpha\_05} is superseded by the document that made it. \\
\addlinespace
2026-08-01 (ruling R6) &
``Reproduce Wu, Zhang \& Zhang's $\beta$ on the $\omega$-destruction term'' &
DAFoam exposes \code{betaFIOmega\_} on the $\omega$ equation's \textbf{production} term
(\code{DAkOmegaSST.C:743}); the destruction term's $\beta$ is the F1-blended model constant
with no field hook. The two are not equivalent, because the other $\omega$ terms do not
scale with $\beta$. What this stack can run today is a production-term inversion, and it is
not a like-for-like reproduction. \\
\addlinespace
2026-08-04 (W2 deep read) &
One paper cited where there are \textbf{two artifacts}; three claims did not survive
reading both &
(i) ``train only on public CBFS'' is true of the \emph{entry} and false of the
\emph{paper}, which trains on mixed NASA hump + CBFS --- and the hump is a scored case, so
the paper's protocol scored on this benchmark would be in-sample by construction;
(ii) ``zero-shot generalise to duct'' is stated by neither artifact and is \emph{our}
inference from the entry's training plus the leaderboard, sound but ours to carry;
(iii) the rank-2 model is \textbf{SST-QCRC}, carrying an untrained QCR2000 quadratic
constitutive term ($c_r = 0.3$) that no laboratory document had noticed. \\
\bottomrule
\end{tabular}
\end{table}

The same practice runs through the rest of the record: a partial withdrawal of the sentence
``no alternative training-family model exists for \code{NASA\_2DWMH}'' once CBFS was found
shipped as legal training data
(\artifact{closure\_challenge\_C6\_hump\_decision.md}, commit
\code{1f5e2403}); the retirement of the ``5 of 8'' headline; and a citation audit that
found two of our own quotations had gone stale \emph{because we improved the pages they
cited} (\S\ref{sec:integrity}).

\subsection{The ducts: what is structurally missing}

The deficit is the ducts. Rank 2 reaches 0.0455 and 0.0399 on \code{AR\_1} and \code{AR\_3}
where we reach 0.0811 and 0.0775; those two cases were 62.9\% of the round-2 gap and remain
the dominant term. Four measured findings explain why, and none of them is fixable by more
model capacity on the present inputs.

\begin{enumerate}
\item \textbf{A proven algebraic degeneracy.} Uncorrected RANS produces \emph{exactly} zero
  secondary flow on every duct: the streamwise fraction of mean $|\mathbf{U}|$ is 1.0000 on
  all four training cases, and the maximum transverse $|\mathbf{U}|$ is $\sim$1e-15 against
  $O(10)$ streamwise, because a linear Boussinesq closure has no mechanism for Prandtl's
  secondary flow of the second kind. For a mean field of that form, \code{I3\_S3} and
  \code{I4\_W2S} vanish identically --- verified on the run's own data at machine epsilon
  ($\sim$1e-14) and independently over 50{,}000 random shear-rate pairs. Two of the seven
  features carry \emph{zero} information on the entire duct family, by mathematical
  necessity.
\item \textbf{And it is worse than we published.} Two further identities were measured and
  had never been recorded: \code{I2\_W2} $= -$\code{I1\_S2} \emph{bit-exactly} (residual 0.0
  on all eight duct cases) and \code{I5\_W2S2} $= -$\code{I1\_S2}$^2/2$ to $3.5\times
  10^{-16}$. \textbf{Five of the seven features carry one independent degree of freedom on
  a duct, not five}; the effective input dimension is 3, or 4 once Variant D's $d/d_{\max}$
  is added, where the physics is three-dimensional. The degeneracy is a property of the
  \emph{baseline we are correcting}, not of duct geometry --- the DNS field does carry
  secondary flow. Control: on the four periodic-hill test cases the same residuals run from
  $2.5\times10^{-3}$ to $9.1\times10^{-1}$, so the test is not vacuous.
\item \textbf{A hypothesis of our own, tested and killed.} Every feature is an $O(3)$
  invariant of $\nabla\mathbf{U}$ or a wall distance, all invariant under the $y
  \leftrightarrow z$ swap that maps the square duct's quadrant onto itself --- which looked
  like a hard error floor built in by construction. Measured on \code{AR\_1\_Ret\_180}, the
  unreachable component is \textbf{0.012\% of the correction's energy}, because the true
  field respects the same symmetry the features do. Refuted, and shipped refuted.
\item \textbf{What binds instead is transfer.} A $k$-NN conditional-variance bound on the
  four training ducts puts the reachable error at 0.011--0.027 against floors of
  0.069--0.107: in-family, the feature set is not the constraint. The binding constraint is
  the \code{Re\_y} extrapolation of \S3.1, and \textbf{no legal test of Reynolds transfer at
  a doubled velocity scale exists inside this benchmark} --- which is why Variant D is
  validated for no-harm on every held-out set the rules permit and is \emph{not} proven to
  close the deficit.
\end{enumerate}

Separately, and stated as a \emph{bound} rather than a decomposition: the missing secondary
flow is 2.07--2.22\% of bulk velocity in the DNS and identically absent from RANS, so it can
explain at most its own magnitude of our duct error --- at most 24\% on \code{AR\_1} and
\code{AR\_3}. \textbf{Scaled MAE does not decompose exactly, so this is a bound and not an
error budget.} Its consequence is nonetheless directional: at least 76\% of our duct error
is streamwise-profile error, and rank 2's \emph{total} duct error (4.55\%, 3.99\%) is
already smaller than our streamwise-only remainder ($\sim$7.0\%, $\sim$6.6\%). A scalar
eddy-viscosity-style correction --- exactly what the field-inversion line targets --- is
well aimed at the dominant term, not wasted against an anisotropy-only gap.

\subsection{What the two entries above us do that we do not}

\paragraph{Versus Reissmann, Fang \& Sandberg (0.0595, gap 0.0059).} Gene-expression-programming
symbolic regression of an anisotropic Reynolds-stress correction, fitted concurrently with
a differential-evolution refit of the $k$--$\omega$ SST coefficients
(lineage: Weatheritt \& Sandberg, \emph{J.\ Comput.\ Phys.} 325 (2016) 22--37;
implementation Reissmann, Fang, Ooi \& Sandberg, \emph{GPEM} 26(1):12 (2025), DOI
10.1007/s10710-025-09510-z). An algebraic stress correction inside the PDE \emph{can
generate} Prandtl's secondary flow of the second kind and carries it across Reynolds number
through the re-solve --- exactly the two axes where they beat us worst. They train on the
same duct family we do, AR 1--10 at $Re_\tau = 180$: \textbf{the difference is not data, it
is model form.} Where we beat them: both \code{alpha\_15} cases (0.0501/0.1011 against
0.0592/0.1339) and both declined \code{alpha\_05} cases.

\paragraph{Versus Wu \& Zhang (0.0624, gap 0.0030).} One training flow, a five-constant
symbolic law, and a re-solve beat our 21-case gradient-boosted model everywhere outside the
periodic-hill family. Their SST-QCRC combines a fixed, untrained QCR-style quadratic
Reynolds-stress term ($c_r = 0.3$, ``directly adopted'', no data-driven training) with a
learned multiplicative field $\beta_{\mathrm{CND}}$ on the $\omega$-destruction term
obtained by conditioned field inversion and PySR symbolic regression; the final law is three
local features collapsed to $\beta_{\mathrm{CND}} = \max(-0.1157,\,0.0058525\,\lambda_2)\cdot
\lambda_2$, bounded below 4 by hand for stability (Wu \& Zhang, \emph{The Training of the
SST-QCRC Model}, the entry's own description document; backing paper Wu, Zhang \& Zhang,
\emph{AIAA J.} 63(2):687--706 (2025), DOI 10.2514/1.J064416, arXiv:2402.16355). What they do
that we do not is \textbf{correct causes rather than symptoms} --- the $\omega$ equation
rather than the velocity field --- and buy generalisation from compactness. The QCR term
supplies duct secondary flow they never trained for.

\paragraph{An external finding worth recording.} Per their own paper's main-text Table 1,
the NASA wall-mounted hump --- the same flow as the scored \code{NASA\_2DWMH} --- is one of
the nine training flows of the rank-3 entry's method (Liu, Wang, Zhao \& Xiao,
arXiv:2509.17189). \textbf{Stated as an observation from their own table, not as an
accusation}: the paper is not the submission, and whether the leaderboard entry retrained
excluding the hump is not determinable from public sources.

%=====================================================================
\section{The field-inversion line: blocked, unblocked, verified}
%=====================================================================

The decomposition points at a scalar eddy-viscosity-style correction inside the equations.
That is the FIML family, and reproducing it is the laboratory's Ladder-B line. Over the
last week it moved from a hard blocker to the first converged, verified closure-relevant
gradient.

\subsection{Capability established, and a deviation disclosed}

Stage 1 (\artifact{dafoam/ladder-b/S1\_FIML\_FIELD\_INVERSION.md},
commit \code{1cd44c04}) established the capability on DAFoam v5.0.0 / OpenFOAM v2506
(He, Mader, Martins \& Maki, \emph{AIAA J.} (2020), \code{J058853}). Four results, in the
order the evidence forces:

\begin{enumerate}
\item \textbf{The mesh-warp defect that fails this laboratory's shape-optimisation cases
  cannot reach a $\beta$ gradient} --- measured, not reasoned. A counter monkeypatched onto
  \code{DAFoamWarper.\_\_init\_\_} and \code{compute\_jacvec\_product} printed
  \code{\{"warper\_init": 0, "warper\_jacvec": 0\}} on every field-inversion adjoint that
  ran to completion. A class never instantiated cannot have a method called on it: the warp
  is not in the model. FIML is reachable on a code path where shape optimisation here is
  not.
\item \textbf{No custom turbulence library is needed}, and a previous record's central
  disclosure was retracted on that basis: \code{betaFIOmega\_} and \code{betaFIK\_} are
  members of DAFoam's own \code{DAkOmegaSST}, constructed \code{READ\_IF\_PRESENT} with
  default 1.0.
\item \textbf{The disclosed deviation.} DAFoam's $\beta$ multiplies the $\omega$ equation's
  \textbf{production} term; the destruction term's $\beta$ is the F1-blended model constant
  and carries no field hook. Anything built on this hook is a production-term inversion and
  must be recorded as such. \textbf{This is a real reduction in what Stage 1 can claim, and
  it is the honest one: building the production version and calling it the roadmap item
  would have been the quieter and worse choice.}
\item \textbf{The adjoint works and is finite-difference verified} on a 5{,}000-cell case
  with 5{,}000 design variables: 2.674\% against the real objective gradient direction at
  step $10^{-3}$ and 2.683\% at $10^{-4}$, with 2.2--3.1\% per resolvable component and
  zero sign flips among components the finite difference can resolve. The residual
  $\sim$2.7\% is \emph{step-independent} (a factor of ten in step moves it by 0.009
  percentage points, where genuine central-difference truncation would fall by
  $\sim$100$\times$) and \emph{convergence-independent} (tightening the primal 1000-fold,
  to an achieved residual of 9.947e-12, moved it from 2.6739\% to 2.6728\%). It is real,
  unexplained, and recorded rather than buried.
\end{enumerate}

A finite-difference \code{check\_totals} over a per-cell $\beta$ field would need $2N$
primal solves --- 10{,}000 on the tutorial case and 103{,}252 on the hump. Field-inversion
gradients cannot be verified the way shape gradients are; the directional-plus-sampled-component
protocol is the affordable substitute, and that is what ``FD-verified'' has to mean at this
design-variable count.

\paragraph{A protocol defect found by auditing our own logs.} A monitor replay flagged four
of 34 finite-difference runs printing DAFoam's velocity-bound message. The escalation asked
which four points to exclude. \textbf{Both the question and its premise were wrong.} DAFoam
gates that message on \code{printInterval}, which the campaign ran at the default of 100 ---
so each log reported clipping from 1\% of its iterations. Re-running twelve points at
\code{printInterval 1} found \textbf{689 clip events against the 3 the archive recorded: a
230$\times$ undercount}, with eleven of twelve runs clipping including the unperturbed
baseline. There was no clip-free subset to retreat to. What the re-run also found is that
all twelve objectives return \emph{bit-identical to sixteen significant figures} with
iteration counts matching exactly, and all 689 events fall in iterations 27--142 of runs
1815--3291 iterations long. The clip is a startup transient the iteration discards; the
verification stands on a narrower and better-stated claim than before --- not ``no clip
occurred'', which was never true and was never measured, but ``the clip is a transient and
the objective is provably unchanged by it''. The real defect found was in the protocol: the
analyser read only the objective line and never inspected the log it came from.

\subsection{The blocker, and its mechanism}

Both closure-relevant target cases failed identically. The DAFoam/PETSc discrete-adjoint
GMRES solve returned \code{PetscConvergedReason: -9} (\code{KSP\_DIVERGED\_NANORINF}) at
iteration 0 on CBFS (21{,}000 cells) and on the NASA hump (51{,}626 cells) --- two meshes,
two objectives, one signature. Memory was excluded by measurement (host \code{MemAvailable}
bottomed at 18.99\,GB with the Jacobian colouring completing normally), as were mesh
quality, wall treatment, a non-differentiable velocity clip, primal convergence tolerance,
and the objective function.

The mechanism was then located rather than guessed
(\artifact{S1\_FIML\_FIELD\_INVERSION.md} \S4 addendum; \code{PROOF.md} \S25.2--25.3):

\begin{itemize}
\item \textbf{The whole matrix-ordering axis was run.} Of DAFoam's five orderings, two
  (\code{rcm}, \code{1wd}) produce the NaN and three (\code{natural}, \code{nd},
  \code{qmd}) produce honest stagnation at \code{-3} with the residual flat to 13
  significant figures; \emph{none} converges. The initial residual is identical
  (7.091590452305e-04) in all five, as it must be. A single buggy ordering would not do
  that; a factorisation hitting a destructive zero pivot under some elimination orders and
  not others does exactly that.
\item \textbf{The cause was reproduced with no DAFoam in the loop.} The preconditioner
  matrix and right-hand side were dumped with stock PETSc runtime flags and analysed
  offline. The matrix (210{,}592 square, 13{,}710{,}468 nonzeros) has no zero rows, columns
  or diagonal entries. \code{scipy.sparse.linalg.spilu} returns \emph{``Factor is exactly
  singular''} --- the \code{-9}, in a second independent implementation --- while
  \code{splu}, full LU \emph{with partial pivoting}, solves the same system to
  $\lVert Ax-b\rVert/\lVert b\rVert = 2.5\times10^{-12}$. The system is nonsingular and
  consistent. \textbf{Fill adds fill, not pivoting}, which is why raising the ILU fill level
  had never helped.
\end{itemize}

\subsection{The unblock, and what it cost}

\artifact{dafoam/ladder-b/W4\_ADJOINT\_PC\_UNBLOCK.md} (commit
\code{1cd44c04}) closes it. The runtime-options route was tested first and found
\emph{reachable but empty}: \code{sub\_}-prefixed factor options do land in the deployed
factor (a 1e-8 zero-pivot threshold is visible in \code{PCView}), correcting a standing
belief that nothing reached the sub-preconditioner --- and every one of them still returns
\code{-9}, so the rebuild was justified by measurement rather than convenience.

The fix is one word. DAFoam hard-codes \code{PCType localPCType = PCILU} for the ASM
sub-blocks; a single-hunk patch makes it \code{PCLU} when
\code{DAFOAM\_SUBPC\_TYPE=lu} is set in the environment, and changes nothing otherwise.
Results:

\begin{itemize}
\item \textbf{CBFS converges} --- \code{PetscConvergedReason: 2}, 667 iterations, 243\,s at
  4 ranks --- on B3's exact \code{-9} configuration, changed only by the environment
  variable. The residual history is instructive: the first 400 iterations crawl, because
  the assembled matrix the LU inverts is not the matrix-free operator GMRES applies, and
  then the residual falls six decades in 267 iterations. \textbf{Slow-then-superlinear is
  what a correct strong preconditioner on a mismatched operator looks like; flat to
  thirteen digits is what the dead ILU looked like. They are not the same thing.}
\item \textbf{The first closure-relevant field-inversion gradient exists}: 21{,}000
  per-cell $\beta$ components on CBFS, objective \code{varianceU} $=
  1.5279278906359758\times10^{-2}$, gradient norm 1.4558e-05.
\item \textbf{The regression control was run before the treatment}: the patched image with
  the switch \emph{unset} reproduces the \code{-9} at iteration 0 with all 13 residual
  digits identical. The rebuild is behaviour-neutral by default.
\item \textbf{On the hump the same switch moves the failure from catastrophic to honest}:
  the factorisation completes and the residual descends monotonically
  (1.094 $\rightarrow$ 0.954 over 900 iterations, accelerating) but does not converge, and
  the run was stopped deliberately when host \code{MemAvailable} fell to 1.62\,GB on a
  shared box. \textbf{The hump blocker is no longer a singular factorisation --- it is
  Krylov convergence rate against a memory envelope}, which is a different, honestly named
  problem.
\end{itemize}

Finite-difference verification of the CBFS gradient met its pre-declared gate at 13--45
times the required tightness (Table~\ref{tab:fd}), with the unperturbed control run
reproducing the adjoint run's objective to \emph{exactly zero difference} across all 17
digits. Total cost, billed at full occupancy for full wall-clock including a 600\,s ordering
timeout: 185.0 core-min against a 240 core-min budget.

\begin{table}[h]
\centering
\small
\caption{Finite-difference verification of the CBFS $\beta$ gradient (21{,}000 design
variables). The last row was produced by an independent verification agent at a cell the
laboratory never published.}
\label{tab:fd}
\begin{tabular}{lccccl}
\toprule
Cell & $h$ & Central FD & Adjoint $g[i]$ & Rel.\ error & Source \\
\midrule
5491 (largest $|g|$) & 0.05 & 1.914384790951e-06 & 1.916018813330e-06 & \textbf{0.085\%} & W4 \S5d \\
5491                 & 0.10 & 1.907246786675e-06 & 1.916018813330e-06 & 0.460\% & W4 \S5d \\
6740                 & 0.05 & 1.679642395377e-06 & 1.678653382471e-06 & \textbf{0.059\%} & W4 \S5d \\
12486                & 0.05 & 1.472401682783e-06 & 1.469472928907e-06 & \textbf{0.199\%} & W4 \S5d \\
6490                 & 0.05 & 1.520027166028e-06 & 1.519705840939e-06 & \textbf{0.021\%} & supervisor sweep \\
\bottomrule
\end{tabular}
\end{table}

The doubling of the error from $h = 0.05$ to $h = 0.1$ on cell 5491, with the finite
difference moving \emph{away} from the adjoint value as the step grows, is
central-difference truncation behaving exactly as $O(h^2)$ predicts: the agreement is
step-limited, not noise-limited, which is the signature of a correct derivative measured
against a resolvable objective.

\subsection{Independent adversarial verification}

The claim was then attacked by a separate verification agent working under the doctrine
that the claim is wrong until it survives attack
(\artifact{dafoam/VERIFICATION\_cbfs\_unblock\_supervisor\_sweep.md},
commit \code{9c19ccc8}). Nothing was taken on the record's word:

\begin{itemize}
\item The patch was diffed \emph{against the built image}, not against the patch file. An
  md5 sweep of every \code{.C}/\code{.H} in the DAFoam source tree of both images shows
  exactly two changed files --- \code{DALinearEqn.C} and its \code{lnInclude} copy ---
  and all three AD-mode libraries contain the new environment-variable string, so all three
  were rebuilt from patched source.
\item The finite-difference table was recomputed from the raw perturbation vectors: each
  perturbation file was verified to differ from the all-ones baseline at \emph{exactly one
  index}, by exactly $\pm h$.
\item A new finite-difference point was derived at cell 6490 --- rank 10 of $|g|$, a cell
  the laboratory never chose --- in an isolated case copy with the sweep's own driver:
  0.0211\%, tighter than any published cell.
\item The env-off regression was re-run cold, independently, reproducing \code{-9} with all
  13 residual digits.
\item The cost ledger was reconciled against log timestamps to the second (693.04\,s
  measured against 693\,s billed; 642\,s summing exactly).
\end{itemize}

Verdict: no defect found, five of five attacks confirmed, one disclosed caveat restated
rather than allowed to go quiet --- the hump attempt of W4 \S5b computes its objective
against \code{NASA\_2DWMH} experimental data, a scored case; it produced no gradient,
nothing from it feeds any entry, and the gate claim rests on CBFS alone.

\subsection{What the rank-2 method actually is, and what a reproduction must copy}

The W2 reading
(\artifact{campaign/W2\_WU\_ZHANG\_DESTRUCTION\_FIML\_READING.md},
commit \code{393ad74f}) read both Wu artifacts in full and returned with the correction of
Table~\ref{tab:corrections} plus an implementation specification. Its operational
consequences are three:

\begin{enumerate}
\item \textbf{A reproduction must copy the \emph{entry's} protocol, not the paper's.} The
  paper's models train on mixed NASA hump and CBFS inversion data; the hump is a scored
  case, so the paper's protocol scored on this benchmark is in-sample by construction. The
  entry's CBFS-only choice is what keeps a reproduction clean.
\item \textbf{A $\beta$-on-destruction patch alone reproduces the paper's inversion
  capability, not the deployed rank-2 model.} Duct parity additionally needs the untrained
  QCR2000 constitutive term, filed separately as
  \code{w3-qcr-constitutive-term-for-rank2-parity}. Any score from a
  $\beta$-destruction-only model must not be set beside 0.0455/0.0399 as like-for-like.
\item \textbf{The patch was specified from the installed source, not from documentation.}
  The $\omega$ equation is assembled in two places (\code{DAkOmegaSST.C:745} in
  \code{calcResiduals()} and \code{:870} in \code{getFvMatrixFields()}) and both carry the
  identical destruction line; patching only one would leave the preconditioner assembling a
  different equation than the residual --- on a stack whose standing blocker \emph{was} a
  factorisation of exactly that matrix, an inconsistent preconditioner would be a new
  confound indistinguishable from the old one. Seven named risks are recorded with the
  specification, including that the conditioned shielding function uses the constant 8 from
  the paper where OpenFOAM's built-in DDES shield would silently supply 20.
\end{enumerate}

\subsection{What is running now}

The Stage 1 CBFS inversion is pre-registered and under way
(\artifact{dafoam/ladder-b/S1\_CBFS\_INVERSION\_PREREGISTRATION.md},
commit \code{e6321e95}), written before any inversion iteration ran and while the only
existing numbers were the baseline objective and the single verified gradient at $\beta = 1$.
It fixes, in advance: the loss
$J(\beta) = \lambda_{\mathrm{QoI}}\,\mathrm{varianceU}(\beta) + \lambda_{L2}\sum_j(\beta_j-1)^2$
with $\lambda_{\mathrm{QoI}} = 65.448114804931393$ so that the data term equals exactly 1 at
$\beta = 1$ (Wu \& Zhang's own normalisation) and $\lambda_{L2} = 10^{-4}$ declared at the
top of their stated band and \emph{not} tuned; bounds $[0.2, 4.0]$, the upper bound theirs
and the lower bound ours by choice and disclosed as such; L-BFGS-B rather than their SLSQP,
because SciPy's SLSQP would allocate a dense $3.749\times10^9$-double workspace --- 30.0\,GB,
the whole of this host's RAM --- at 21{,}000 design variables; a hard budget cap of 600
core-min with a stop rule that records a capped run as \emph{capped, not converged}; and two
success criteria, a 30\% reduction in the data term and a requirement that more than half of
the top-decile $\beta$ deviations fall inside a pre-declared spatial window around the
recirculation bubble and separated shear layer.

It also fixes its own label before any result exists: \textbf{this is a production-term
inversion}, at $\sim$36 evaluations against Wu \& Zhang's $\sim$140 SLSQP iterations, with a
field-variance loss rather than their 30 sparse points --- three independent reasons no
number it produces may be set beside theirs as like-for-like.

%=====================================================================
\section{Verification and integrity}
\label{sec:integrity}
%=====================================================================

\subsection{Leakage guards, in code rather than in prose}

The one strict rule was audited line by line in the code rather than taken from our own
documentation (\artifact{CLOSURE\_CHALLENGE\_SUBMISSION\_DRAFT.md} \S4.1):

\begin{itemize}
\item Ground-truth reads occur only inside loops over the 21 training cases. The test-case
  loop calls the RANS-field loader only, annotated \code{\# no U\_LES read}.
\item The split is enforced by executable assertions, not comments:
  \code{assert len(\_PH\_TRAIN) == 21 and len(\_PH\_VAL) == 4}, and assertions that the
  training, validation and test name sets do not intersect. The test list is imported
  \emph{solely} to assert non-intersection.
\item The CSV exporter does not merely omit scoring calls: before any pipeline work it
  replaces \code{score}, \code{score\_from\_csv}, \code{evaluate\_by\_case},
  \code{evaluate\_from\_csv\_by\_case}, \code{evaluate\_individual\_case} and the
  ground-truth accessors with raising stubs across three namespaces, then \emph{proves the
  guard is armed} by calling \code{score()} and catching the refusal before continuing. If
  a future edit reintroduces a scoring call, the script dies rather than quietly spending
  one.
\item An in-sample gate (\code{closure\_in\_sample\_gate.py}) runs before and after
  closure work and passes.
\end{itemize}

\paragraph{The residual soft leakage, disclosed by us before anyone asks.} Across rounds,
test feedback influenced \emph{which machinery got built}. The round-3 record states it in
its own words: \emph{``The gate's MOTIVATION came from round 2's official call \dots\ That
is soft, adaptive leakage, and it is real.''} The defence is that the gate's every parameter
derives from 21 training cases and was reproduced byte-for-byte from train-only data before
touching test features; that seeing round-2's scores changed \emph{when} it was applied, not
\emph{what} it outputs, so its four test decisions would have been identical had round 2
never been scored; and that the benchmark itself ships the ground truth and instructs
submitters to preview per-case scores, so an interpretation under which observing them is a
violation would make the benchmark's own documented workflow a violation. \textbf{Verdict:
not a violation of the rule as written; residual risk low but non-zero, and interpretive
rather than factual.} It belongs in the submission's cover material so that a reviewer
\emph{finds} it written down rather than discovers it.

\subsection{The scoring-call ledger}
\label{sec:ledger}

Five distinct prediction sets have ever been scored against the eight test cases' ground
truth: the RANS-identity floor, and rounds 1 through 4. The unit counted is
\emph{distinct prediction sets}, which is a defensible and meaningful unit --- and the
compliance audit caught a docstring claiming something stronger and demonstrably false
(``exactly one call'' where the same run makes four invocations across two prediction sets,
the extra pair re-scoring the already-counted floor as a harness check). The claim was
substantively harmless and was fixed anyway, because \textbf{an entry whose entire
credibility rests on precise self-accounting cannot ship with a sentence a reviewer can
falsify by reading forty lines further down the same file.}

The ledger is a self-imposed discipline. The benchmark imposes no limit and instructs
submitters to preview their score. No public statement of ours may imply otherwise. Nor
may any statement of ours imply the ledger is principled: in Blum \& Hardt's own words
(arXiv:1502.04585, \S2.3 above) rate limiting is a ``poorly understood heuristic'', and
naming ours as one is the honest price of it.

\subsection{The adversarial methods audit}
\label{sec:audit}

An audit conducted on 2026-08-04 proof-checked the closure line end to end under the
instruction to assume the records overstate until verified, test-blind throughout
(\artifact{CLOSURE\_METHODS\_COMPARISON.md}, commit \code{ac2f37ee}). The metric definition
held as written. Six gaps were found in our own evidence, ranked by how much the entry's
credibility rests on the unverified assumption, each costed as a proposal:

\begin{description}
\item[G1 --- seed sensitivity, real and live.] \code{random\_state=0} everywhere and no
  model ever retrained under another seed. This is not pedantry: at 327{,}600 training
  cells the PH model crosses scikit-learn 1.9.0's 200{,}000-row binning subsample, so the
  bin edges --- hence the trees, hence three of the eight predictions --- depend on the
  seed, against a gap to rank 2 of 0.0030.
\item[G2 --- no physicality check on any corrected field.] $\nabla\cdot\mathbf{U}$ never
  measured, on the one dimension every other entrant gets by construction.
\item[G3 --- baseline convergence evidence exists but is uncited] for all eight
  predictions. It matters twice over on the two declined cases, where the shipped solve
  \emph{is} our submission. \textbf{Closed 2026-08-05}
  (\artifact{CLOSURE\_CHALLENGE\_STATUS.md} \S4b, commit \code{11ff63b1}): the four
  periodic-hill test cases carry
  final residuals of $\sim$1e-10 after 20{,}000 iterations, and a sweep of all 29 cases
  reproduces the pattern on 28 of 29 --- with the exception recorded honestly rather than
  smoothed, one \emph{training} case whose shipped residual file is truncated at iteration
  76, so that its convergence is plausible but not evidenced. The three duct test cases
  terminate with \emph{``SIMPLE solution converged''} at 405 / 1540 / 7009 iterations under
  explicit \code{residualControl}, and the $O(0.4)$--$O(0.7)$ transverse residuals that
  look alarming are recorded as a normalisation artifact on a machine-zero field
  ($\max|U_y| = 2.1\times10^{-15}$ against $U_x$ up to 108\,m/s), noted so that nobody
  later mistakes them for a divergent solve.
\item[G4 --- manifest gap.] The three new round-4 duct CSVs carried no recorded SHA-256 and
  the round-4 submission directory had no manifest; the hashing discipline covered only the
  five copied files. \textbf{Closed 2026-08-05}
  (\artifact{closure\_challenge\_submission\_round4/MANIFEST.json}, commit \code{30fd6d7a}):
  all eight files hashed, zero scoring calls, guaranteed by a script that never imports the
  benchmark package at all.
\item[G5 --- gate statistical power] is thin and already at its disclosed limit; widen from
  4 held-out decisions to 25 by leave-one-out over all periodic-hill cases.
\item[G6 --- no training-convergence curve] for either gradient-boosted model.
\end{description}

None requires a scoring call; none touches test ground truth; all six keep the leakage
guards armed. Four are now closed: G1 and G2 were taken first, on pre-registered thresholds
(\S\ref{sec:stability}), and G3 and G4 landed on 2026-08-05.

\subsection{Seed stability and field physicality: the lines drawn before the measurements}
\label{sec:stability}

The pre-registration for G1 and G2 was committed
(\artifact{closure\_challenge\_stability\_physicality\_audit.md} \S0,
commit \code{73fa6a33}; results at \code{fee5eaf9} and \code{cbe040f6}) \emph{before} either
run existed, so the materiality lines cannot have been drawn around the results:

\begin{itemize}
\item \textbf{G1:} retrain the PH model at eight seeds with everything else frozen, score on
  the four validation cases only, and compute a test-side prediction-spread bound requiring
  no ground truth at all. \emph{Material if} the overall-equivalent induced spread reaches
  0.0003 (10\% of the gap to rank 2) by either estimate, \emph{or} if any single case's
  spread reaches 0.0030 --- a one-case spread the size of the whole gap being material
  regardless of the mean.
\item \textbf{G2:} push both the raw RANS and the corrected field through the
  \emph{identical} already-validated Green--Gauss operator so that discretisation error
  cancels in the comparison, and report the volume-weighted RMS divergence and its ratio.
  \emph{A ratio of 2 or more on any model-corrected case is recorded as a material
  physicality cost of the method class and must be carried into the submission's disclosure
  section; a ratio below 2 is still reported per case, quantified, with no editorialising
  either way.}
\end{itemize}

\paragraph{G1, measured.} The results are recorded in \S1 of the same audit file (commit
\code{fee5eaf9}), with raw numbers in
\artifact{closure\_challenge\_seed\_sensitivity.json} (generated
2026-08-04T18:55Z; zero scoring calls; guard re-armed and proven armed first; ground truth
read for the 21 training and 4 validation cases only). Seed 0 reproduces the recorded
round-1 model exactly (validation pooled 0.0876) and its predictions match the shipped
round-4 CSVs to 5e-10. Across seeds 0--7 the pooled validation metric spans only 0.000276,
and the two proxy ranges for the served \code{alpha\_15} predictions are 0.000568 and
0.000520; the overall-equivalent proxy estimate is $S_{\text{proxy}} = 0.000261$, just below
the 0.0003 line. But the test-side spread \emph{bounds} --- the largest disagreement between
any two seeds' own predictions at the 1000 evaluation points, normalised the same way the
metric is --- reach \textbf{0.0085} and \textbf{0.0099} on the two \code{alpha\_15} cases
and 0.0010 on the hump, giving $S_{\text{bound}} = 0.00242$.

\textbf{Verdict: material}, by two of the three pre-declared triggers --- $S_{\text{bound}}
\geq 0.0003$, and two single-case bounds at or above 0.0030 --- and both halves of the
reading belong in the record:

\begin{itemize}
\item \textbf{Where truth is legitimately available, the metric is stable under seed.} On
  the four held-out validation cases the per-case spread is 0.0005--0.0012 and the
  overall-equivalent is $\sim$0.00026, a tenth of the gap to rank 2. That is the best
  available \emph{estimate} of the seed effect on the score, and it is small: the entry
  gains a measured rather than assumed stability property, an estimated one-seed
  uncertainty on the overall of about 0.0003.
\item \textbf{Where truth is not available, the only honest statement is the bound, and it
  is not small.} Across seeds the trees move the pointwise test predictions by 0.85--0.99\%
  of the velocity scale on the two served \code{alpha\_15} cases --- per-case movement up to
  roughly three times the entire gap \emph{if it failed to cancel against the truth}. On
  validation cases that wobble demonstrably cancels to the 0.0006 level in the metric;
  expecting the same on test is reasonable, but it is an extrapolation and not a
  measurement, and the pre-registered rule counts the bound. Tightening it would cost a
  scoring call and is not taken.
\end{itemize}

The consequence was fixed in advance and has been executed: \textbf{our leaderboard-gap
language now carries a seed-uncertainty qualifier} in the status record and on the public
page. Scope is exact --- the duct model behind three predictions is bin-deterministic at
41{,}971 cells, so \code{random\_state} is inert for it by construction rather than by
measurement, and the two declined cases are the organisers' own solve, so the five non-PH
predictions contribute zero seed variance and the division by eight is exact.

\paragraph{G2, measured --- and the finding is against us.} The divergence audit completed
on 2026-08-05 (310\,s on the two-core cap, zero scoring calls, 37 cases; results in \S2 of
the audit file, raw record
\artifact{closure\_challenge\_divergence\_audit.json}). Every
regenerated test field was matched against its shipped round-4 CSV \emph{before} its
divergence number was believed --- maximum absolute difference 4.8e-10 to 4.9e-8, all at CSV
write precision --- so the numbers are about the submitted fields and not a lookalike.

\begin{table}[h]
\centering
\small
\caption{Volume-weighted RMS divergence of the eight submitted fields, both fields through
the identical validated Green--Gauss operator. The scale-free column is the continuity error
relative to the field's own velocity-gradient magnitude, $\lVert\nabla\cdot\mathbf{U}\rVert/\lVert\nabla\mathbf{U}\rVert_F$, given for the RANS field and for the submitted field; the RANS value is the
operator-consistency floor, since the solver enforces continuity on face fluxes rather than
on reconstructed cell gradients.}
\label{tab:div}
\resizebox{\textwidth}{!}{%
\begin{tabular}{llccc}
\toprule
Case & Submitted field & Ratio & rel.\ RANS & rel.\ corrected \\
\midrule
\code{alpha\_15\_13929\_4048} & PH-corrected      & \textbf{58.0}  & 0.18\% & \textbf{10.5\%} \\
\code{alpha\_15\_13929\_2024} & PH-corrected      & \textbf{123.6} & 0.08\% & \textbf{9.7\%} \\
\code{alpha\_05\_4071\_4048}  & declined (raw RANS) & 1.000        & 0.47\% & 0.47\% \\
\code{alpha\_05\_4071\_2024}  & declined (raw RANS) & 1.000        & 0.27\% & 0.27\% \\
\code{AR\_1\_Ret\_360}        & duct-D corrected  & $\sim$1e15\,$^{\dagger}$ & 2.4e-17 & \textbf{2.3\%} \\
\code{AR\_3\_Ret\_360}        & duct-D corrected  & $\sim$4e14\,$^{\dagger}$ & 6.9e-17 & \textbf{3.1\%} \\
\code{AR\_14\_Ret\_180}       & duct-D corrected  & $\sim$7e14\,$^{\dagger}$ & 4.6e-17 & \textbf{3.4\%} \\
\code{NASA\_2DWMH}            & PH-corrected      & 1.526          & 0.43\% & 0.66\% \\
\bottomrule
\end{tabular}}

\smallskip
{\footnotesize $^{\dagger}$ The duct ratios are divisions by machine zero and should be read
as such: a fully developed, exactly unidirectional RANS duct field is cell-wise
divergence-free to roundoff (9e-12 against a gradient scale of $4\times10^5$). The
interpretable statement is the last column, not the ratio.}
\end{table}

\textbf{Verdict: material.} Thirty-five of thirty-six model-corrected cases sit at or above
the pre-registered ratio of 2, and five of the six corrected submissions exceed it; only
\code{NASA\_2DWMH} (1.53) stays under, for the consistent reason that the PH-trained
correction is small relative to that case's own gradient scale --- the same model that
barely moved its score barely moves its continuity error. The train and validation cases
show the same corrected-field levels (8--19\% on hills, 2--3\% on ducts), so this is a
property of the method on its own turf and not a test-set artifact.

Three things follow, and none of them is an excuse:

\begin{itemize}
\item \textbf{The cost is real and now has a number.} The correction takes fields whose
  measured continuity error sits at the operator floor and returns fields violating
  continuity at $\sim$10\% of the gradient scale on the two corrected hills and 2.3--3.4\%
  on the three ducts. That is far above the operator floor, so it is the field and not the
  operator.
\item \textbf{The scoring metric never sees it.} Nothing here changes 0.0654. But it is
  exactly what a referee of the description document will ask, and per the pre-registered
  consequence it is carried into the submission's disclosure section --- ``we measured it
  and here is the number'' beats silence, even when the number is against us.
\item \textbf{The two declined cases are, again, the strongest rows.} Their ratio is exactly
  1.000 because the gate's ``off'' state ships the organisers' own solve and inherits its
  physicality untouched. The one part of our submission that is physical by construction is
  the part where our model does not appear.
\end{itemize}

This is also the measured case for the field-inversion line of \S5. The argument for moving
the correction \emph{inside} the equations had been an argument from taxonomy and from the
duct decomposition; it is now an argument from a number on our own submitted fields. A
correction that re-enters the PDE and is re-solved returns a field satisfying continuity
because the solver enforces it --- which is precisely what the four entries above and below
us get for free, and precisely what Stage 1 is being built to give us.

\subsection{The pre-registration that said no}
\label{sec:nogo}

The clearest evidence that the discipline binds is the round-5 candidate
(\artifact{closure\_challenge\_R5\_ALPHA05\_REGIME\_PREREGISTRATION.md},
commit \code{be642415}). C2 had recommended training a second, regime-specific model as
strictly better than gating, because it recovers the term instead of merely declining to
lose it. The candidate was built --- 11 of 21 training cases qualified under a membership
rule defined by the same frozen C1 gate that makes the test-time decision, and the gate
refit reproduced $\alpha = 0.7499$ and threshold 0.1263 before membership was computed, the
script exiting rather than defining membership from a drifted gate.

The go/no-go rules were declared in the committed training script's docstring before any
leave-one-out delta existed. Measured:

\begin{table}[h]
\centering
\small
\caption{The R5 pre-declared gate, and what the measurements did to it.}
\begin{tabular}{L{6.2cm}L{6.6cm}c}
\toprule
Rule & Measured & Verdict \\
\midrule
(i) the validation case improves &
\code{alpha\_05\_10071\_4048}: 0.0759 $\rightarrow$ \textbf{0.0447} ($-$0.0312), where the
global round-1 model gives 0.0772, a hurt & PASS \\
(ii) mean leave-one-out delta $< 0$ &
$-$0.0195 over 10 cases, 8 of 10 improved, two by more than 0.035 & PASS \\
(iii) no case hurt by more than $+$0.010 &
\code{alpha\_05\_4071\_3036}: 0.0492 $\rightarrow$ 0.0659, \textbf{$+$0.0167} & \textbf{FAIL} \\
\bottomrule
\end{tabular}
\end{table}

\textbf{Verdict: NO-GO. The rules were not revisited after the numbers, the scoring call was
not authorised, and nothing was scored.} Round 4 stands as the entry of record.

\paragraph{The rule has a name, and ours is the weaker version of it (added 2026-08-05).}
``Adopt the retrained model only if it is certified not to be worse than the incumbent''
is a named research problem: Thomas, Theocharous \& Ghavamzadeh, \emph{High Confidence
Policy Improvement}, \emph{Proc.\ 32nd ICML}, PMLR 37 (2015) 2380--2388 (abstract read
2026-08-05 from the PMLR record). Their algorithm gives ``probabilistic guarantees about
the quality of each policy that it proposes'', with the user free to ``select any
performance lower-bound and confidence level.'' Two differences cut against us and one
for us. Against: their bound is probabilistic and user-parameterised, where
\textbf{our $+$0.010 per-case hurt cap is a single hand-chosen number whose
pre-registration never states a derivation} --- it was declared before the deltas
existed, which is the part that matters for leakage, but ``why 0.010'' is unanswered in
our own record, and the method-priority review
(\artifact{CLOSURE\_METHOD\_PRIORITY\_REVIEW.md} \S2.4, \S7, commit \code{d84b649f})
files that as a charter trigger with a proposal,
\code{agenda/proposals/w8-a-hurt-cap-that-states-where-it-came-from.json}. Also against:
clinical trials pre-specify a \textbf{non-inferiority margin} in the protocol for exactly
the reason we froze ours, and our cap is that idea under another name rather than bespoke
machinery. For us: their guarantee is about the aggregate, ours is per-case, which is
stricter --- and it is why R5 failed on one case while improving on eight of ten.

The finding underneath is worth more than the candidate would have been (recorded as lesson
L-33). The case that fails is \code{alpha\_05\_4071\_3036} --- the training sibling of both
declined test cases, sharing their $\alpha$ and their 4071 parameter and differing only in
the third index --- and it is the lowest-baseline case in the entire training split
(0.0492). The leave-one-out table splits cleanly on baseline quality: every case with
baseline $\geq 0.058$ improves and both cases below $\sim$0.055 get worse, in both variants.
\textbf{C2's mechanism is recursive: train only on the low-baseline regime and the
hurt-where-RANS-is-already-good boundary does not disappear --- it moves down, from
$\sim$0.07--0.13 to $\sim$0.055, and reappears inside the regime.} A correction model fitted
to a population still damages the best-baseline members of that population.

A successor candidate --- a second-stage threshold declining even the regime model below
$\sim$0.055 --- is obvious and is explicitly \emph{not} adopted, because it was conceived
after those numbers existed and adopting it now would be exactly the
rules-revisited-after-the-fact move the document's own \S3 forbids. It is filed as a
recommendation needing its own pre-registration, its own membership audit, and a held-out
design that does not lean on a single in-regime validation case.

\subsection{Corrections and withdrawals register}

\begin{itemize}
\item \textbf{The published duct diagnosis was wrong} and was corrected by the round that
  acted on it (\S3.1).
\item \textbf{``No alternative training-family model exists for \code{NASA\_2DWMH}'' was
  withdrawn}: CBFS --- the curved backward-facing step, the nearest flow in this benchmark
  to a wall-mounted hump --- ships as \emph{training} data with its LES fields and is legal
  to fit on. The pre-registered argument that chose the periodic-hill model was a two-way
  comparison of hills against ducts in which CBFS is not named, not rejected and not
  mentioned: the choice was defended against the worse of two candidates while the better
  third sat unused in the training set. What still keeps the case shut is not the absence of
  a candidate but \emph{who} may choose and \emph{when} --- anyone choosing today has seen
  round 2's score, so taking that route legally requires a pre-registration frozen first.
  Both routes are priced; neither has been taken.
\item \textbf{``We are last on the ducts'' was never true.} A working belief repeated in
  internal notes was re-derived from the harness and found false; after round 4 we are last
  on \code{NASA\_2DWMH} alone.
\item \textbf{Two of our own quotations went stale because we improved the pages they
  cited.} A citation audit extracted all 28 quotations in the submission draft mechanically
  and matched them against normalised source text. Both defects were in citations of sources
  \emph{we control} --- not one was in a citation of the benchmark's rules, which had been
  transcribed carefully because they were known to matter. The audit closes the question
  that no sentence in that document is attributed to a source nobody here has read, and
  explicitly does \emph{not} close the general problem: nothing checks staleness
  automatically, so it must be re-run at the moment of approval, not before.
\item \textbf{An over-generalisation in our own solver record was corrected on
  measurement}: a prior finding that a certain PETSc failure code fools DAFoam's success
  gate was found to be specific to \code{-5}; \code{-9} never fools it. It matters in the
  safe direction --- the blocked cases were never at risk of silently returning a wrong
  gradient.
\end{itemize}

%=====================================================================
\section{Outlook}
%=====================================================================

\subsection{The technical path}

\begin{description}
\item[Duct parity is the largest available term.] Matching the best published entry on the
  two aspect ratios we trail is worth about 0.011 on the overall --- more than the entire
  0.0059 gap to rank 1. Every move currently visible in that family is either contaminated
  by knowledge of per-case test outcomes or needs evidence this benchmark cannot supply, so
  a structural route has been filed instead: the scalar-invariant feature set is provably
  three-dimensional on a duct while the tensor basis is not. Measured on two duct training
  cases, the second basis tensor carries 1.87--2.36$\times$ the norm of the first, with 71\%
  of its magnitude and 91\% of the third tensor's in the transverse block that generates
  secondary flow. That is a specific structural argument for the duct family belonging to a
  tensor-basis ansatz.
\item[Rank-2 parity needs two components, not one.] The $\beta$-on-destruction model patch
  is specified line by line against installed source and filed; the untrained QCR
  constitutive term is filed separately. Landing only the first reproduces the paper's
  inversion capability, not the deployed rank-2 model, and no score from it may be compared
  like-for-like.
\item[Stage 1 is running.] The CBFS production-term inversion is pre-registered, budgeted at
  600 core-min, checkpointed every evaluation, and labelled in advance. It can establish that
  the verified gradient drives an actual descent on a closure-relevant case, where a
  production-term $\beta$ wants to move on CBFS, and the real per-iteration price that Stage
  2 planning needs. It cannot establish anything about destruction-term physics, about
  generalisation, or about any benchmark score --- and it says so before it runs.
\item[The hump remains open, honestly named.] Its blocker is now Krylov convergence rate
  against a memory envelope rather than a singular factorisation. A Richardson-wrapped
  configuration capping the restart basis is staged and unrun, awaiting a fresh budget.
\item[Audit items.] Four of the six are closed: G1 and G2 both returned material findings
  against pre-registered thresholds, and G3 and G4 landed on 2026-08-05. What remains is
  G5, which widens the gate's evidence base from 4 held-out decisions to 25 by
  leave-one-out over all periodic-hill cases, and G6, which is free if folded into the
  retraining runs.
\end{description}

\subsection{Submission status}

The package exists and is audited. The eight CSVs are written, 1000 rows by 3 columns,
comma-delimited and header-free, verified finite and round-trip accurate below 5e-8, with
mean velocity magnitudes falling inside the band spanned by all four accepted submissions on
all eight cases --- a check that would catch a point-ordering, column-ordering or units
blunder at zero scoring cost. Both defects the compliance audit found (the false docstring;
the non-existent submittable artifact) are closed, the evaluation package is cited by commit
hash rather than by its misleading version string, and the flagged preprint quotations have
been verified verbatim.

The description document's disclosures are fixed and non-negotiable: the method class and
its weakness relative to the family it is ranked against; that two submitted fields are the
uncorrected baseline and any leaderboard-best status there belongs to the baseline; the full
adaptive-leakage history in our own words; the executable train/validation/test assertions;
that the hump regression was not reverted and why; the duct feature degeneracy as a negative
result we would rather publish than have inferred from our scores; the scoring-call count
correctly framed as self-imposed; the measured continuity cost of \S\ref{sec:stability} and
the measured seed-uncertainty qualifier on the rank-2 gap; the prior art, split as \S2.2
splits it into classifiers that \emph{identify} and classifiers that \emph{control}, with
no sentence presenting confidence-gated correction as novel; and --- added 2026-08-05 and
mandatory --- the Hanna et al.\ citation for the method class itself, with its delta
stated (\S\ref{sec:method}).

\medskip
\noindent What remains is not technical. Two items are outstanding and both are Katie's: the
author line and reference URL, and the proofread and approval without which nothing moves.
Two questions are drafted for the steward and unsent --- whether a company name is acceptable
on a leaderboard row, and whether any terms of use apply to a commercial entity training on
data whose repositories carry no licence file at all. \textbf{As of today the only thing
between this package and the steward is Katie.}

\appendix
%=====================================================================
\section{Provenance index}
%=====================================================================

\begin{longtable}{L{7.4cm}L{1.9cm}L{5.4cm}}
\caption{Repository artifacts cited in this report. Paths are relative to
\code{demo-output/website/} unless otherwise shown; commits are those at which the cited
text stands.}\label{tab:prov}\\
\toprule
Artifact & Commit & What it carries \\
\midrule
\endfirsthead
\toprule
Artifact & Commit & What it carries \\
\midrule
\endhead
\artifact{CLOSURE\_CHALLENGE\_STATUS.md} & \code{b9cb56a2} &
Rounds 1--4 history; the metric; per-case tables; feature-expressivity and generalisation
studies \\
\artifact{CLOSURE\_METHODS\_COMPARISON.md} & \code{ac2f37ee} &
Metric read from the evaluation code; entrant dossiers; the dimension-by-dimension
comparison; gap findings G1--G6 \\
\artifact{closure\_challenge\_C2\_error\_decomposition.md} &
\code{1da013a8} & The error decomposition and its corrections \\
\artifact{CLOSURE\_CHALLENGE\_SUBMISSION\_DRAFT.md} & \code{15530f97} &
Policy, eligibility, compliance audit, draft package, addenda \\
\artifact{CLOSURE\_CHALLENGE\_PRIOR\_ART.md} & \code{15530f97} &
Prior art on classifier-gated corrections: \S2.1 identify, \S2.3--\S2.4 control, \S2.7
the PMoE near miss read in full \\
\artifact{CLOSURE\_METHOD\_PRIORITY\_REVIEW.md} & \code{d84b649f} &
Priority and novelty review: 27 searches, 13 counted negatives, the eight owed citations
and the two phrasing corrections applied in \S2.1--\S2.3, \S\ref{sec:ledger} and
\S\ref{sec:nogo} \\
\artifact{closure\_challenge\_C1\_error\_estimator.md} &
\code{5719374e} & The decline gate: fit, validation, and its stated statistical limits \\
\artifact{closure\_challenge\_C6\_hump\_decision.md} & \code{1f5e2403} &
Both hump routes priced; neither taken \\
\artifact{closure\_challenge\_R5\_ALPHA05\_REGIME\_PREREGISTRATION.md} &
\code{be642415} & The round-5 candidate and its NO-GO \\
\artifact{closure\_challenge\_stability\_physicality\_audit.md} &
\code{73fa6a33} (\S0), \code{cbe040f6} (results) & G1/G2 thresholds committed before either
run; results \S1--\S2 \\
\artifact{closure\_challenge\_seed\_sensitivity.json} & \code{cbe040f6}
& G1 raw record: 8 seeds, spreads, bounds, verdict \\
\artifact{closure\_challenge\_divergence\_audit.json} & \code{cbe040f6}
& G2 raw record: 37 cases, RMS divergence before and after correction \\
\artifact{closure\_challenge\_submission\_round4/MANIFEST.json} &
\code{30fd6d7a} & All eight submission CSVs hashed; closes G4 \\
\artifact{closure\_challenge\_trained\_entry\_round4\_duct.json} &
\code{4c2f3562} & The entry of record \\
\artifact{closure\_challenge\_duct\_reynolds\_transfer.json} &
\code{d0992da2} & The round-4 pre-registration \\
\artifact{dafoam/ladder-b/S1\_FIML\_FIELD\_INVERSION.md} &
\code{1cd44c04} & FIML capability, FD verification, the clip audit, the blocker \\
\artifact{dafoam/ladder-b/W4\_ADJOINT\_PC\_UNBLOCK.md} &
\code{1cd44c04} & The sub-block LU unblock and the first closure-relevant gradient \\
\artifact{dafoam/VERIFICATION\_cbfs\_unblock\_supervisor\_sweep.md} &
\code{9c19ccc8} & Independent adversarial verification of that claim \\
\artifact{dafoam/ladder-b/S1\_CBFS\_INVERSION\_PREREGISTRATION.md} &
\code{e6321e95} & The inversion now running, pre-registered \\
\artifact{campaign/W2\_WU\_ZHANG\_DESTRUCTION\_FIML\_READING.md} &
\code{393ad74f} & Both Wu artifacts read in full; the patch specification \\
\bottomrule
\end{longtable}

\section{External works cited}

\begin{itemize}
\item McConkey, Buchanan, Smidt, Bodner, Dwight \& Cinnella. \emph{The Closure Challenge}.
  arXiv:2603.28884, DOI 10.48550/arXiv.2603.28884.
\item Wu, C., Zhang, S. \& Zhang, Y. \emph{Development of a Generalizable Data-driven
  Turbulence Model: Conditioned Field Inversion and Symbolic Regression}. \emph{AIAA
  Journal} 63(2):687--706 (2025), DOI 10.2514/1.J064416, arXiv:2402.16355.
\item Wu, C. \& Zhang, Y. \emph{The Training of the SST-QCRC Model} --- the rank-2 entry's
  own description document, held at
  \artifact{docs/papers/wu\_zhang\_sst\_qcrc\_challenge\_description.pdf}.
\item Reissmann, Fang, Ooi \& Sandberg. \emph{Genetic Programming and Evolvable Machines}
  26(1):12 (2025), DOI 10.1007/s10710-025-09510-z; lineage Weatheritt \& Sandberg,
  \emph{J.\ Comput.\ Phys.} 325 (2016) 22--37.
\item Liu, Wang, Zhao \& Xiao. \emph{Toward a unified data-driven turbulence model through
  multi-objective learning}. arXiv:2509.17189.
\item Oulghelou, Cherroud, Merle \& Cinnella. \emph{Machine-Learning-Assisted Blending of
  Data-Driven Turbulence Models}. \emph{Flow, Turbulence and Combustion} (2025),
  DOI 10.1007/s10494-025-00661-8, arXiv:2410.14431.
\item Hanna, B.\,N., Dinh, N.\,T., Youngblood, R.\,W. \& Bolotnov, I.\,A.
  \emph{Coarse-Grid Computational Fluid Dynamic (CG-CFD) Error Prediction using Machine
  Learning}. arXiv:1710.09105 (2017); journal version \emph{Progress in Nuclear Energy}
  118 (2019) 103140, DOI 10.1016/j.pnucene.2019.103140. \textbf{Read in full}
  (arXiv text, held at
  \artifact{docs/papers/hanna\_dinh\_youngblood\_bolotnov\_1710.09105.pdf}); the journal
  version is cited at metadata-only tier.
\item Kennedy, M.\,C. \& O'Hagan, A. \emph{Bayesian Calibration of Computer Models}.
  \emph{J.\ R.\ Statist.\ Soc.\ B} 63(3) (2001) 425--464. \emph{Metadata only --- not
  read in this laboratory; cited as lineage, nothing quoted.}
\item Wu, J.-L., Xiao, H., Sun, R. \& Wang, Q. \emph{Reynolds-averaged Navier--Stokes
  equations with explicit data-driven Reynolds stress closure can be ill-conditioned}.
  \emph{J.\ Fluid Mech.} 869 (2019) 553--586, arXiv:1803.05581. \emph{Abstract read.}
\item Ji, H., Luo, Y., Zhou, H. \& Zhao, Y. \emph{Progressive Mixture-of-Experts with
  autoencoder routing for continual RANS turbulence modelling}. arXiv:2601.09305 (2026).
  \textbf{Read in full.}
\item Blum, A. \& Hardt, M. \emph{The Ladder: A Reliable Leaderboard for Machine Learning
  Competitions}. arXiv:1502.04585 (2015); ICML 2015, PMLR 37. \emph{Abstract read.}
\item NeurIPS Workshop on Pre-registration in Machine Learning, PMLR 148 (2020) and
  PMLR 181 (2021). \emph{Metadata only.}
\item Thomas, P., Theocharous, G. \& Ghavamzadeh, M. \emph{High Confidence Policy
  Improvement}. \emph{Proc.\ 32nd ICML}, PMLR 37 (2015) 2380--2388. \emph{Abstract read.}
\item Duraisamy, K., Iaccarino, G. \& Xiao, H. \emph{Turbulence Modeling in the Age of
  Data}. \emph{Annu.\ Rev.\ Fluid Mech.} 51 (2019) 357--377. \emph{Metadata only --- not
  read in this laboratory; the ``correct the answer'' phrasing is our own paraphrase via
  \artifact{docs/research/CLOSURE\_METHODS.md}, not a quotation of the review.}
\item Ling, J. \& Templeton, J. \emph{Phys.\ Fluids} 27, 085103 (2015).
\item Wu, Wang, Xiao \& Ling. \emph{Flow, Turbulence and Combustion} 99, 25 (2017),
  DOI 10.1007/s10494-017-9807-0, arXiv:1607.04563.
\item Steiner, Dwight \& Vir\'e. \emph{Flow, Turbulence and Combustion} (2022),
  DOI 10.1007/s10494-022-00346-6, arXiv:2106.15593.
\item Buchanan, L\u{a}c\u{a}tu\c{s}, West \& Dwight. \emph{Computers and Fluids} 305,
  106899 (2025), DOI 10.1016/j.compfluid.2025.106899, arXiv:2504.06758.
\item Schmelzer, Dwight \& Cinnella (SpaRTA). \emph{Flow, Turbulence and Combustion} 104,
  579--603, DOI 10.1007/s10494-019-00089-x.
\item Pope, S.\,B. \emph{A more general effective-viscosity hypothesis}. \emph{J.\ Fluid
  Mech.} 72 (1975) 331--340.
\item He, Mader, Martins \& Maki. \emph{DAFoam: An Open-Source Adjoint Framework for
  Multidisciplinary Design Optimization with OpenFOAM}. \emph{AIAA Journal} (2020),
  \code{J058853}; and \emph{Computers and Fluids} (2018), discrete adjoint for OpenFOAM.
\item Bentaleb, Lardeau \& Leschziner (2012) --- the curved backward-facing step LES that
  supplies the field-inversion reference data, as shipped by the benchmark for its training
  case.
\end{itemize}

\vfill
\noindent\rule{\textwidth}{0.4pt}\\
{\footnotesize Prepared by the Certonomous CFD laboratory, 5 August 2026. Nothing in this
document has been sent to the benchmark steward. Every position claim is dated; the
challenge is ongoing and the public leaderboard can change without notice.}

\end{document}
